\subsection{Essentially integrable functions}

Let F be a real Banach space; recall that the elements of the spaces \(F_{F}^{p}\) (Ch. IV, §3, No.~3) and \(L_{F}^{p}\) (Ch. IV, §3, No.~4, Def. 2) are \(\mu\) -moderated functions (Ch. IV, §5, No.~6, Lemma 1); with \(N_{F}\) still denoting the space of negligible mappings of T into F, we shall introduce the space \(N_{F}^{\infty}\) of locally negligible mappings of T into F.

Lemma. --- Let g and \(g'\) be two \(\mu\) -moderated mappings with values in F; if g and \(g'\) are equal locally almost everywhere to a same function f, then \(g = g'\) almost everywhere.

For, let D be the set of \(t \in T\) such that \(\mathbf{g}(t) \neq \mathbf{g}'(t)\) ; D is locally negligible and moderated, therefore negligible (Cor. 1 of Prop. 7).

We shall denote by \(\overline{\mathcal{F}}_F^p (\mathrm{T},\mu)\) (or simply \(\overline{\mathcal{F}}_{\mathrm{F}}^{p}(\mu),\overline{\mathcal{F}}_{\mathrm{F}}^{p}\) , if no confusion can result) the set of mappings \(\mathbf{f}\) of \(\mathrm{T}\) into \(\mathrm{F}\) , such that there exists a function \(\mathbf{g}\in \mathcal{F}_{\mathrm{F}}^{p}\) equal to \(\mathbf{f}\) locally almost everywhere. Since the number \(\mathrm{N}_p(\mathbf{g})\) depends only on \(\mathbf{f}\) by the Lemma, we will write \(\overline{\mathrm{N}}_p(\mathbf{f}) = \mathrm{N}_p(\mathbf{g})\) . The function \(\overline{\mathrm{N}}_p\) is obviously a semi-norm on \(\overline{\mathcal{F}}_F^p\) , and we shall always assume that \(\overline{\mathcal{F}}_F^p\) is equipped with the topology defined by \(\overline{\mathrm{N}}_p\) . The closure of 0 for this topology is the space \(\mathcal{N}_{\mathrm{F}}^{\infty}\) ; the relations \(\overline{\mathcal{F}}_{\mathrm{F}}^{p} = \mathcal{F}_{\mathrm{F}}^{p} + \mathcal{N}_{\mathrm{F}}^{\infty}\) , \(\mathcal{N}_{\mathrm{F}}^{\infty}\cap \mathcal{F}_{\mathrm{F}}^{p} = \mathcal{N}_{\mathrm{F}}\) (Lemma) show that the normed space \(\overline{\mathcal{F}}_{\mathrm{F}}^{p} / \mathcal{N}_{\mathrm{F}}^{\infty}\) may be canonically identified with \(\mathcal{F}_{\mathrm{F}}^{p} / \mathcal{N}_{\mathrm{F}}\) , which is complete (Ch. IV, §3, No.~3, Prop. 5); therefore \(\overline{\mathcal{F}}_F^p\) is itself complete.

We shall similarly denote by \(\overline{\mathcal{L}}_F^p (\mathrm{T},\mu)\) (or \(\overline{\mathcal{L}}_F^p (\mu)\) , or \(\overline{\mathcal{L}}_F^p\) ) the subspace \(\mathcal{L}_F^p +\mathcal{N}_F^\infty\) of \(\overline{\mathcal{F}}_F^p\) ; one can also characterize \(\overline{\mathcal{L}}_F^p\) as the subspace of \(\overline{\mathcal{F}}_F^p\) constituted by the measurable mappings (Ch. IV, §5, No.~6, Th. 5).

The normed space \(\overline{L}_{F}^{p}/N_{F}^{\infty}\) may be canonically identified with \(L_{F}^{p}\) ; \(\overline{L}_{F}^{p}\) is therefore complete. Its elements are called the p-th power essentially integrable functions, this terminology being justified by the following proposition:

PROPOSITION 9. --- For a mapping f of T into F to belong to \(\overline{\mathcal{F}}_F^p\) (resp. to \(\overline{\mathcal{L}}_F^p\) ), if is necessary and sufficient that (resp. that f be measurable and that)

\[
\mu^ {\bullet} (| \mathbf {f} | ^ {p}) <   + \infty .
\]

One then has \(\overline{\mathrm{N}}_p(\mathbf{f}) = (\mu^\bullet (|\mathbf{f}|^p))^{1 / p}\) .

We may clearly limit ourselves to the assertion concerning \(\overline{F}_{F}^{p}\) . If f belongs to \(\overline{F}_{F}^{p}\) , let g be a function belonging to \(F_{F}^{p}\) that is equal to f locally almost everywhere; then \(|f|^{p}=|g|^{p}\) locally almost everywhere, therefore

\[
\mu^ {\bullet} (| \mathbf {f} | ^ {p}) = \mu^ {\bullet} (| \mathbf {g} | ^ {p}) = \mu^ {*} (| \mathbf {g} | ^ {p}) <   + \infty
\]

(Prop. 1, a) and Prop. 7), and, on the other hand, by the definition of \(\overline{\mathbf{N}}_p\) ,

\[
\overline {{\mathrm{N}}} _ {p} (\mathbf {f}) = \mathrm{N} _ {p} (\mathbf {g}) = \left(\mu^ {*} (| \mathbf {g} | ^ {p})\right) ^ {1 / p}.
\]

Conversely, suppose that \(\mu^{\bullet}(|\mathbf{f}|^{p}) < +\infty\) ; then there exists a moderated set A such that f is zero locally almost everywhere in T-A (Prop. 7). The function \(f\varphi_{A}\) , equal locally almost everywhere to f, is such that \(\mathrm{N}_{p}(\mathbf{f}\varphi_{\mathrm{A}}) = \overline{\mathrm{N}}_{p}(\mathbf{f}) < +\infty\) , therefore it belongs to \(F_{F}^{p}\) , and \(f \in F_{F}^{p}\) .

COROLLARY. --- For f to belong to \(L_{F}^{p}\) , it is necessary and sufficient that f belong to \(\overline{L}_{F}^{p}\) and be moderated.

DEFINITION 3. -- The elements of \(\overline{\mathcal{L}}_F^1\) are called essentially \(\mu\) -integrable functions with values in F. On composing the mapping \(\widetilde{\mathbf{f}} \mapsto \mu(\mathbf{f})\) of \(L sub F superscript 1\) into F with the canonical mapping of \(\overline{\mathcal{L}}_F^1\) onto \(L sub F superscript 1\) , one obtains a continuous linear mapping of \(\overline{\mathcal{L}}_F^1\) into F that extends the mapping \(\mathbf{f} \mapsto \int \mathbf{f} d\mu\) of \(\mathcal{L}_F^1\) into F. One again denotes by \(\int \mathbf{f} d\mu\) or \(\mu(\mathbf{f})\) the value of this mapping for \(\mathbf{f} \in \overline{\mathcal{L}}_F^1\) , and this element is called the integral of \(\mathbf{f}\) with respect to \(\mu\) .

Two essentially integrable functions that are equal locally almost everywhere have the same integral. For every function \(f \geqslant 0\) that is finite and essentially integrable, \(\int^{\bullet} f d\mu = \int f d\mu\) . If A is a set whose characteristic function is essentially integrable, then A is said to be an essentially \(\mu\) -integrable set; \(\int \varphi_{A} d\mu\) is also denoted \(\mu(A)\) and is again called the measure of A.

If a function \(\mathbf{f}\) , with values in \(\mathbf{F}\) , is defined locally almost everywhere in \(\mathbf{T}\) , we again say that \(\mathbf{f}\) is essentially integrable if it is equal, locally almost everywhere, to a function \(\mathbf{f}_1\) that is everywhere defined and integrable; we then set

\[
\int \mathbf {f} d \mu = \int \mathbf {f} _ {1} d \mu ,
\]

and this definition is independent of the integrable function \(f_{1}\) everywhere defined and equal locally almost everywhere to f (Lemma). One defines similarly the notion of essentially integrable function for functions with values in \(\overline{R}\) that are defined and finite locally almost everywhere.

The reader will have no difficulty in extending, to essentially integrable functions, the results of Ch. IV, §4 for integrable functions, on replacing `almost everywhere' in the statements by `locally almost everywhere.' We note for example the inequality

\[
\left| \int \mathbf {f} d \mu \right| \leqslant \int | \mathbf {f} | d \mu ,\tag{3}
\]

valid for every essentially integrable function f with values in a Banach space.

PROPOSITION 10. --- Let \(\mathfrak{K}\) be a \(\mu\) -dense set of compact subsets of T. a) If \(f\) is a numerical function \(\geqslant 0\) , then

\[
\mu^ {\bullet} (f) = \sup _ {K \in \mathfrak {K}} \mu^ {*} (f \varphi_ {K}).\tag{4}
\]

\begin{enumerate}
\def\labelenumi{\alph{enumi})}
\setcounter{enumi}{1}
\tightlist
\item
  If \(\mathbf{f}\) is an essentially integrable function with values in a Banach space \(\mathbf{F}\) , then
\end{enumerate}

\[
\int \mathbf {f} d \mu = \lim _ {\mathfrak {K}} \int \mathbf {f} \varphi_ {\mathrm{K}} d \mu ,
\]

the limit being taken with respect to the directed (for \(\subset\) ) set \(\mathfrak{K}\) .

To establish a), it suffices to show that for every compact subset L of T, \(\int^{*}f\varphi_{L}d\mu=\sup_{K}\int^{*}f\varphi_{K}d\mu\) , where K runs over the set of subsets of L belonging to K. Since L is the union of a negligible set and an increasing sequence \((K_{n})\) of elements of K (Ch. IV, §5, No.~8, Prop. 12), this follows from the theorem on passage to the limit in upper integrals (Ch. IV, §1, No.~3, Th. 3).

Suppose now that f belongs to \(\overline{L}_{F}^{1}\) ; let \(\varepsilon\) be a number \textgreater0, and let K be an element of K such that

\[
\int | \mathbf {f} | \varphi_ {\mathrm{K}} d \mu \geqslant \int | \mathbf {f} | d \mu - \varepsilon
\]

(such a K exists by a)). Then, for every compact set H containing K,

\[
\left| \int \mathbf {f} d \mu - \int \mathbf {f} \varphi_ {\mathrm{H}} d \mu \right| \leqslant \int | \mathbf {f} | \varphi_ {\mathbf {C H}} d \mu \leqslant \int | \mathbf {f} | \varphi_ {\mathbf {C K}} d \mu \leqslant \varepsilon .
\]

Extension to complex Banach spaces and measures. Let F be a complex Banach space; by an abuse of notation, the real Banach space underlying F will also be denoted by F. The Banach space \(\overline{\mathcal{L}}_{\mathrm{F}}^{p}(\mathrm{T},\mu)\) may then be equipped with a natural complex Banach space structure, and it is necessary to be specific as to whether one is using the real or the complex structure of this space. In this chapter, and absent express mention to the contrary, it will always be understood to be the real structure.

Let \(\theta\) be a complex measure; we set \(\overline{\mathcal{L}}_F^p (\mathrm{T},\theta) = \overline{\mathcal{L}}_F^p (\mathrm{T},|\theta |)\) ; if \(\mathbf{F}\) is a complex Banach space, one can make the same remarks as above. In particular, a function \(\mathbf{f}\) with values in \(\mathbf{F}\) will be called essentially integrable for \(\theta\) if it is essentially integrable for \(|\theta |\) . Assertion b) of Prop. 10 then extends at once to complex measures.

